% Vorgehen und Aufbau der Arbeit (Kap. 1.3, fig:aufbau).
% Eigene Darstellung, keine Quelle. Kapiteltitel wie in Kapitel/*.tex (Stand 25.11.2026), gekürzt.
% Stil: ba-grafik.tex. Einbinden: % Vorgehen und Aufbau der Arbeit (Kap. 1.3, fig:aufbau).
% Eigene Darstellung, keine Quelle. Kapiteltitel wie in Kapitel/*.tex (Stand 25.11.2026), gekürzt.
% Stil: ba-grafik.tex. Einbinden: % Vorgehen und Aufbau der Arbeit (Kap. 1.3, fig:aufbau).
% Eigene Darstellung, keine Quelle. Kapiteltitel wie in Kapitel/*.tex (Stand 25.11.2026), gekürzt.
% Stil: ba-grafik.tex. Einbinden: % Vorgehen und Aufbau der Arbeit (Kap. 1.3, fig:aufbau).
% Eigene Darstellung, keine Quelle. Kapiteltitel wie in Kapitel/*.tex (Stand 25.11.2026), gekürzt.
% Stil: ba-grafik.tex. Einbinden: \input{Bilder/tikz/aufbau-arbeit} in einer figure-Umgebung. Nicht skalieren.
% Raster: vier Spalten, Mitten x = 1,7 | 5,7 | 9,7 | 13,7, Blockbreite 34 mm (Kanten 0..15,4), Blockhöhe 14 mm,
% Zeilenabstand 1,8. Treppe 5 -> 6 -> 7 -> 8 -> 9 = Herleitung, Entwurf, Umsetzung, Prüfung, Zusammenfassung.
% Zeilenumbrüche von Hand, ohne text width, damit in den Kästen keine Silbentrennung entsteht. Geprüft im 12pt-Build.
\begin{tikzpicture}[ba]
  \tikzset{
    k/.style={ba block, minimum width=34mm, minimum height=14mm, font=\footnotesize},
    kh/.style={k, draw=baGruen, fill=baFlaeche},
    kopf/.style={font=\small\bfseries, anchor=south},
    klein/.style={font=\footnotesize, text=baGrau, align=center, anchor=north},
  }
  % Spaltenköpfe
  \node[kopf] at (1.7,0.95)  {Grundlagen};
  \node[kopf] at (5.7,0.95)  {Handlungsbedarf};
  \node[kopf] at (9.7,0.95)  {Entwicklung};
  \node[kopf] at (13.7,0.95) {Bewertung};
  % Kapitel
  \node[k]  (K2) at (1.7,0)     {2 Humanoide Roboter\\und Unitree G1};
  \node[k]  (K3) at (1.7,-1.8)  {3 Simulation und\\Infrastruktur};
  \node[k]  (K4) at (1.7,-3.6)  {4 Fähigkeiten,\\verwandte Arbeiten};
  \node[kh] (K5) at (5.7,0)     {5 Handlungsbedarf,\\Anforderungen,\\Methodik};
  \node[kh] (K6) at (9.7,0)     {6 Konzept};
  \node[kh] (K7) at (9.7,-1.8)  {7 Umsetzung};
  \node[kh] (K8) at (13.7,-1.8) {8 Evaluation};
  \node[k]  (K9) at (13.7,-3.6) {9 Zusammenfassung\\und Ausblick};
  % Grundlagen bündeln: Klammer rechts der Spalte, ein Pfeil zu Kap. 5
  \draw[ba normal] (K2.east) -- (3.75,0);
  \draw[ba normal] (K3.east) -- (3.75,-1.8);
  \draw[ba normal] (K4.east) -- (3.75,-3.6) -- (3.75,0);
  \draw[ba pfeil] (3.75,0) -- (K5.west);
  % Fluss
  \draw[ba pfeil] (K5.east) -- (K6.west);
  \draw[ba pfeil] (K6.south) -- (K7.north);
  \draw[ba pfeil] (K7.east) -- (K8.west);
  \draw[ba pfeil] (K8.south) -- (K9.north);
  % Rolle der Spalte für die Forschungsfragen
  \node[klein] at (1.7,-4.55)  {Stand der Technik,\\wertungsfrei};
  \node[klein] at (5.7,-0.95)  {Defizite, F1 bis F4,\\Anforderungen,\\Versuchsplan};
  \node[klein] at (9.7,-2.75)  {Lösung zu den\\Anforderungen};
  \node[klein] at (13.7,-4.55) {Antworten auf\\F1 bis F4};
\end{tikzpicture}
 in einer figure-Umgebung. Nicht skalieren.
% Raster: vier Spalten, Mitten x = 1,7 | 5,7 | 9,7 | 13,7, Blockbreite 34 mm (Kanten 0..15,4), Blockhöhe 14 mm,
% Zeilenabstand 1,8. Treppe 5 -> 6 -> 7 -> 8 -> 9 = Herleitung, Entwurf, Umsetzung, Prüfung, Zusammenfassung.
% Zeilenumbrüche von Hand, ohne text width, damit in den Kästen keine Silbentrennung entsteht. Geprüft im 12pt-Build.
\begin{tikzpicture}[ba]
  \tikzset{
    k/.style={ba block, minimum width=34mm, minimum height=14mm, font=\footnotesize},
    kh/.style={k, draw=baGruen, fill=baFlaeche},
    kopf/.style={font=\small\bfseries, anchor=south},
    klein/.style={font=\footnotesize, text=baGrau, align=center, anchor=north},
  }
  % Spaltenköpfe
  \node[kopf] at (1.7,0.95)  {Grundlagen};
  \node[kopf] at (5.7,0.95)  {Handlungsbedarf};
  \node[kopf] at (9.7,0.95)  {Entwicklung};
  \node[kopf] at (13.7,0.95) {Bewertung};
  % Kapitel
  \node[k]  (K2) at (1.7,0)     {2 Humanoide Roboter\\und Unitree G1};
  \node[k]  (K3) at (1.7,-1.8)  {3 Simulation und\\Infrastruktur};
  \node[k]  (K4) at (1.7,-3.6)  {4 Fähigkeiten,\\verwandte Arbeiten};
  \node[kh] (K5) at (5.7,0)     {5 Handlungsbedarf,\\Anforderungen,\\Methodik};
  \node[kh] (K6) at (9.7,0)     {6 Konzept};
  \node[kh] (K7) at (9.7,-1.8)  {7 Umsetzung};
  \node[kh] (K8) at (13.7,-1.8) {8 Evaluation};
  \node[k]  (K9) at (13.7,-3.6) {9 Zusammenfassung\\und Ausblick};
  % Grundlagen bündeln: Klammer rechts der Spalte, ein Pfeil zu Kap. 5
  \draw[ba normal] (K2.east) -- (3.75,0);
  \draw[ba normal] (K3.east) -- (3.75,-1.8);
  \draw[ba normal] (K4.east) -- (3.75,-3.6) -- (3.75,0);
  \draw[ba pfeil] (3.75,0) -- (K5.west);
  % Fluss
  \draw[ba pfeil] (K5.east) -- (K6.west);
  \draw[ba pfeil] (K6.south) -- (K7.north);
  \draw[ba pfeil] (K7.east) -- (K8.west);
  \draw[ba pfeil] (K8.south) -- (K9.north);
  % Rolle der Spalte für die Forschungsfragen
  \node[klein] at (1.7,-4.55)  {Stand der Technik,\\wertungsfrei};
  \node[klein] at (5.7,-0.95)  {Defizite, F1 bis F4,\\Anforderungen,\\Versuchsplan};
  \node[klein] at (9.7,-2.75)  {Lösung zu den\\Anforderungen};
  \node[klein] at (13.7,-4.55) {Antworten auf\\F1 bis F4};
\end{tikzpicture}
 in einer figure-Umgebung. Nicht skalieren.
% Raster: vier Spalten, Mitten x = 1,7 | 5,7 | 9,7 | 13,7, Blockbreite 34 mm (Kanten 0..15,4), Blockhöhe 14 mm,
% Zeilenabstand 1,8. Treppe 5 -> 6 -> 7 -> 8 -> 9 = Herleitung, Entwurf, Umsetzung, Prüfung, Zusammenfassung.
% Zeilenumbrüche von Hand, ohne text width, damit in den Kästen keine Silbentrennung entsteht. Geprüft im 12pt-Build.
\begin{tikzpicture}[ba]
  \tikzset{
    k/.style={ba block, minimum width=34mm, minimum height=14mm, font=\footnotesize},
    kh/.style={k, draw=baGruen, fill=baFlaeche},
    kopf/.style={font=\small\bfseries, anchor=south},
    klein/.style={font=\footnotesize, text=baGrau, align=center, anchor=north},
  }
  % Spaltenköpfe
  \node[kopf] at (1.7,0.95)  {Grundlagen};
  \node[kopf] at (5.7,0.95)  {Handlungsbedarf};
  \node[kopf] at (9.7,0.95)  {Entwicklung};
  \node[kopf] at (13.7,0.95) {Bewertung};
  % Kapitel
  \node[k]  (K2) at (1.7,0)     {2 Humanoide Roboter\\und Unitree G1};
  \node[k]  (K3) at (1.7,-1.8)  {3 Simulation und\\Infrastruktur};
  \node[k]  (K4) at (1.7,-3.6)  {4 Fähigkeiten,\\verwandte Arbeiten};
  \node[kh] (K5) at (5.7,0)     {5 Handlungsbedarf,\\Anforderungen,\\Methodik};
  \node[kh] (K6) at (9.7,0)     {6 Konzept};
  \node[kh] (K7) at (9.7,-1.8)  {7 Umsetzung};
  \node[kh] (K8) at (13.7,-1.8) {8 Evaluation};
  \node[k]  (K9) at (13.7,-3.6) {9 Zusammenfassung\\und Ausblick};
  % Grundlagen bündeln: Klammer rechts der Spalte, ein Pfeil zu Kap. 5
  \draw[ba normal] (K2.east) -- (3.75,0);
  \draw[ba normal] (K3.east) -- (3.75,-1.8);
  \draw[ba normal] (K4.east) -- (3.75,-3.6) -- (3.75,0);
  \draw[ba pfeil] (3.75,0) -- (K5.west);
  % Fluss
  \draw[ba pfeil] (K5.east) -- (K6.west);
  \draw[ba pfeil] (K6.south) -- (K7.north);
  \draw[ba pfeil] (K7.east) -- (K8.west);
  \draw[ba pfeil] (K8.south) -- (K9.north);
  % Rolle der Spalte für die Forschungsfragen
  \node[klein] at (1.7,-4.55)  {Stand der Technik,\\wertungsfrei};
  \node[klein] at (5.7,-0.95)  {Defizite, F1 bis F4,\\Anforderungen,\\Versuchsplan};
  \node[klein] at (9.7,-2.75)  {Lösung zu den\\Anforderungen};
  \node[klein] at (13.7,-4.55) {Antworten auf\\F1 bis F4};
\end{tikzpicture}
 in einer figure-Umgebung. Nicht skalieren.
% Raster: vier Spalten, Mitten x = 1,7 | 5,7 | 9,7 | 13,7, Blockbreite 34 mm (Kanten 0..15,4), Blockhöhe 14 mm,
% Zeilenabstand 1,8. Treppe 5 -> 6 -> 7 -> 8 -> 9 = Herleitung, Entwurf, Umsetzung, Prüfung, Zusammenfassung.
% Zeilenumbrüche von Hand, ohne text width, damit in den Kästen keine Silbentrennung entsteht. Geprüft im 12pt-Build.
\begin{tikzpicture}[ba]
  \tikzset{
    k/.style={ba block, minimum width=34mm, minimum height=14mm, font=\footnotesize},
    kh/.style={k, draw=baGruen, fill=baFlaeche},
    kopf/.style={font=\small\bfseries, anchor=south},
    klein/.style={font=\footnotesize, text=baGrau, align=center, anchor=north},
  }
  % Spaltenköpfe
  \node[kopf] at (1.7,0.95)  {Grundlagen};
  \node[kopf] at (5.7,0.95)  {Handlungsbedarf};
  \node[kopf] at (9.7,0.95)  {Entwicklung};
  \node[kopf] at (13.7,0.95) {Bewertung};
  % Kapitel
  \node[k]  (K2) at (1.7,0)     {2 Humanoide Roboter\\und Unitree G1};
  \node[k]  (K3) at (1.7,-1.8)  {3 Simulation und\\Infrastruktur};
  \node[k]  (K4) at (1.7,-3.6)  {4 Fähigkeiten,\\verwandte Arbeiten};
  \node[kh] (K5) at (5.7,0)     {5 Handlungsbedarf,\\Anforderungen,\\Methodik};
  \node[kh] (K6) at (9.7,0)     {6 Konzept};
  \node[kh] (K7) at (9.7,-1.8)  {7 Umsetzung};
  \node[kh] (K8) at (13.7,-1.8) {8 Evaluation};
  \node[k]  (K9) at (13.7,-3.6) {9 Zusammenfassung\\und Ausblick};
  % Grundlagen bündeln: Klammer rechts der Spalte, ein Pfeil zu Kap. 5
  \draw[ba normal] (K2.east) -- (3.75,0);
  \draw[ba normal] (K3.east) -- (3.75,-1.8);
  \draw[ba normal] (K4.east) -- (3.75,-3.6) -- (3.75,0);
  \draw[ba pfeil] (3.75,0) -- (K5.west);
  % Fluss
  \draw[ba pfeil] (K5.east) -- (K6.west);
  \draw[ba pfeil] (K6.south) -- (K7.north);
  \draw[ba pfeil] (K7.east) -- (K8.west);
  \draw[ba pfeil] (K8.south) -- (K9.north);
  % Rolle der Spalte für die Forschungsfragen
  \node[klein] at (1.7,-4.55)  {Stand der Technik,\\wertungsfrei};
  \node[klein] at (5.7,-0.95)  {Defizite, F1 bis F4,\\Anforderungen,\\Versuchsplan};
  \node[klein] at (9.7,-2.75)  {Lösung zu den\\Anforderungen};
  \node[klein] at (13.7,-4.55) {Antworten auf\\F1 bis F4};
\end{tikzpicture}
